% ---------------------------------------------------------------------------
%  Staff AI Engineer Preparation Kit -- preamble
%
%  All machinery lives here. chapters/ carry content and nothing else, so a
%  chapter file can be read as prose and a change to the look of the book
%  touches one file.
%
%  The three claim macros at the bottom are the part worth understanding
%  before editing anything: they are what stops this book stating opinion as
%  fact, and they are written so the wrong form is a build error rather than
%  something a reviewer has to notice.
% ---------------------------------------------------------------------------

\usepackage[T1]{fontenc}
\usepackage[utf8]{inputenc}

% Optional packages, probed rather than required.
%
% This is not politeness. The machine this book was written on carries a bare
% TeX Live with neither lmodern nor inconsolata; CI carries a fuller one. A
% preamble that demanded either would build in exactly one of those two places
% and the failure would arrive as somebody else's problem. Probing means the
% book sets in both, and the only difference is which fonts it sets in --
% which moves line breaks and therefore page counts. Do not quote a page count
% from one installation as a fact about the other.
\IfFileExists{lmodern.sty}{\usepackage{lmodern}}{}
\IfFileExists{inconsolata.sty}{\usepackage[varqu]{inconsolata}}{%
  \IfFileExists{courier.sty}{\usepackage{courier}}{}}
% expansion=false is not a preference. Font expansion needs a scalable font,
% and a TeX Live without lmodern falls back to bitmap Computer Modern, where
% pdfTeX refuses outright: "auto expansion is only possible with scalable
% fonts", fatal, no PDF. Protrusion alone works everywhere and is most of the
% benefit. Turning expansion on would build on the fuller installation and die
% on the barer one, which is the same defect as any other unprobed dependency.
\IfFileExists{microtype.sty}{\usepackage[expansion=false]{microtype}}{}

\usepackage[a4paper,
            top=25mm, bottom=28mm,
            inner=30mm, outer=25mm,
            headsep=8mm, footskip=14mm,
            head=15pt]{geometry}

\usepackage{xcolor}
\usepackage{graphicx}
\usepackage{booktabs}
\usepackage{tabularx}
\usepackage{xltabular}
\usepackage{environ}
\usepackage{enumitem}
\usepackage{needspace}
\usepackage[most]{tcolorbox}
\usepackage{titlesec}
\usepackage{fancyhdr}
\usepackage{amsmath}
\usepackage{url}

% Figures are drawn in TikZ and compiled with the book.
%
% The estate's other knowledge bases render diagrams from committed Mermaid
% sources through a browser and a package registry, and a check keeps the two
% copies byte-identical. That is the right shape for a book with three hundred
% diagrams. This one has four, and the pipeline would make them buildable only
% where npm resolves and a headless browser starts -- so a reader cloning this
% repository could not rebuild the book. Recorded as D-9 in DEVIATIONS.md with
% the condition that would reverse it.
\usepackage{tikz}
\usetikzlibrary{positioning, arrows.meta, fit, backgrounds, calc}

\tikzset{
  kitbox/.style={
    draw=kitink, fill=white, line width=0.6pt, rounded corners=1pt,
    inner sep=4pt, align=center, font=\small},
  kitsoft/.style={
    draw=kitgrey, fill=kitgrey!8, line width=0.5pt, rounded corners=1pt,
    inner sep=4pt, align=center, font=\small},
  kitwarnbox/.style={
    draw=kitwarn, fill=white, line width=0.7pt, rounded corners=1pt,
    inner sep=4pt, align=center, font=\small},
  kitarrow/.style={-{Stealth[length=4pt]}, draw=kitink, line width=0.6pt},
  kitlabel/.style={font=\scriptsize\itshape, text=kitgrey},
}

% One figure environment, so every figure in the book has the same caption
% treatment and the same placement rules.
\newcommand{\kitfig}[3]{%
  \begin{figure}[htbp]
  \centering
  #2
  \caption{#3}
  \label{fig:#1}
  \end{figure}}

% upquote makes ' an ASCII apostrophe whatever the typewriter font is. Without
% it a straight quote in a code fragment sets as a right single quotation mark,
% and the line stops being code you can paste. It is font-dependent, which is
% why it ships as a defect rather than as a visible mistake.
\IfFileExists{upquote.sty}{\usepackage{upquote}}{}

\usepackage[hidelinks]{hyperref}
\hypersetup{
  colorlinks=true,
  linkcolor=kitink,
  urlcolor=kitlink,
  citecolor=kitink,
  pdftitle={Staff AI Engineer Preparation Kit},
  pdfsubject={Preparing for a Staff-level AI engineering hiring process},
  pdfkeywords={ai engineering, staff engineer, interview preparation, evals},
  pdflang={en-GB}
}

% ---------------------------------------------------------------------------
%  Palette
%
%  Four colours and no more. A fifth would mean the first four are not
%  carrying meaning.
% ---------------------------------------------------------------------------
\definecolor{kitink}{RGB}{31,78,121}     % structure: headings, rules, links
\definecolor{kitlink}{RGB}{31,78,121}
\definecolor{kitsaid}{RGB}{21,101,64}    % what the market reports
\definecolor{kitmine}{RGB}{140,70,20}    % what the author thinks
\definecolor{kitwarn}{RGB}{150,30,30}    % a trap
\definecolor{kitgrey}{RGB}{110,110,110}

% ---------------------------------------------------------------------------
%  Headings and running heads
%
%  \raggedright on chapter titles is load-bearing rather than cosmetic: at
%  this size a justified title on this measure has nowhere to break and sets
%  an overfull line that no amount of rewording reliably clears.
% ---------------------------------------------------------------------------
\titleformat{\chapter}[display]
  {\normalfont\huge\bfseries\color{kitink}\raggedright}
  {\normalsize\normalfont\color{kitgrey}\MakeUppercase{\chaptertitlename}\ \thechapter}
  {8pt}{\Huge\raggedright}
  [\vspace{2pt}{\color{kitink}\rule{\linewidth}{1.2pt}}]

\titleformat{\section}
  {\normalfont\Large\bfseries\color{kitink}\raggedright}{\thesection}{0.7em}{}
\titleformat{\subsection}
  {\normalfont\large\bfseries\raggedright}{\thesubsection}{0.6em}{}

\titlespacing*{\section}{0pt}{2.2ex plus 1ex minus .2ex}{1.1ex plus .2ex}
\titlespacing*{\subsection}{0pt}{1.8ex plus 1ex minus .2ex}{0.8ex plus .2ex}

\pagestyle{fancy}
\fancyhf{}
\fancyhead[LE,RO]{\small\thepage}
\fancyhead[RE]{\small\itshape\nouppercase{\leftmark}}
\fancyhead[LO]{\small\itshape\nouppercase{\kitheadcap{\rightmark}}}
\renewcommand{\headrulewidth}{0.4pt}
\renewcommand{\footrulewidth}{0pt}

% An over-long running head does not overfull. fancyhdr sets an over-wide
% field at its natural width and prints it on top of the field beside it --
% no warning, no error, and the page number underneath somebody's chapter
% title. Measure the box and scale only when it is too wide, reserving room
% for the folio so the reserve tracks the point size.
\newlength{\kitheadroom}
\newcommand{\kitheadcap}[1]{%
  \setlength{\kitheadroom}{\dimexpr\headwidth-4em\relax}%
  \sbox0{\small\itshape #1}%
  \ifdim\wd0>\kitheadroom
    \resizebox{\kitheadroom}{!}{\usebox{0}}%
  \else
    \usebox{0}%
  \fi}

\renewcommand{\chaptermark}[1]{\markboth{#1}{}}
\renewcommand{\sectionmark}[1]{\markright{#1}}

% A last resort, and a deliberately small one.
%
% This book is set at 12pt on a 155mm measure and it carries long technical
% compounds and file names. Without any emergency stretch TeX produces a tail
% of small overfull lines that no rewording reliably clears, because the
% offending word is usually a name. 1em is enough to absorb that tail and small
% enough that a genuinely unbreakable run still reports.
%
% The risk this carries is real and worth naming: emergency stretch hides a
% wide box by loosening the line instead of fixing it. Which is why it is 1em
% rather than 3em, and why the build still reports what is left.
\setlength{\emergencystretch}{1em}

\setlength{\parskip}{0pt}
\setlength{\parindent}{1.2em}
\setlist{itemsep=2pt, topsep=4pt, parsep=0pt}

% ---------------------------------------------------------------------------
%  Boxes -- two treatments, and that is the whole visual grammar
%
%  A tinted block with a thick left bar is something the reader may take in
%  passing. A full rule on white is something they must stop at. A third
%  treatment would mean the first two are not carrying their meaning.
%
%  before skip and after skip are the numbers that matter. tcolorbox defaults
%  to a couple of points, which is not enough air to read as a separate block,
%  and every box carrying its own value is how a page ends up looking
%  arbitrary. One \tcbset, eight boxes, one set of spacing numbers.
% ---------------------------------------------------------------------------
\tcbset{
  kit tinted/.style={
    enhanced, breakable, sharp corners,
    colback=#1!6!white, colframe=#1!6!white,
    borderline west={2.5pt}{0pt}{#1},
    boxrule=0pt, left=8pt, right=8pt, top=6pt, bottom=6pt,
    before skip=11pt, after skip=11pt,
    fonttitle=\bfseries\small, coltitle=#1,
    attach boxed title to top left={xshift=6pt, yshift=-2pt},
    boxed title style={colback=white, colframe=white, sharp corners, boxrule=0pt}
  },
  kit outlined/.style={
    enhanced, breakable, sharp corners,
    colback=white, colframe=#1,
    boxrule=0.9pt, left=8pt, right=8pt, top=6pt, bottom=6pt,
    before skip=11pt, after skip=11pt,
    fonttitle=\bfseries\small, coltitle=#1,
    attach boxed title to top left={xshift=6pt, yshift*=-2pt},
    boxed title style={colback=white, colframe=white, sharp corners, boxrule=0pt},
    overlay first and middle={%
      \node[anchor=south east, inner sep=2pt, font=\scriptsize\color{#1}]
        at (frame.south east) {$\triangleright$};}
  }
}

\newtcolorbox{kitnote}[1][]{kit tinted=kitink, title={Note}, #1}
\newtcolorbox{kittrap}[1][]{kit outlined=kitwarn, title={The trap}, #1}
\newtcolorbox{kitmethod}[1][]{kit outlined=kitink, title={The method}, #1}
\newtcolorbox{kitrigour}[1][]{kit tinted=kitgrey, title={What this rests on}, #1}

% ---------------------------------------------------------------------------
%  THE THREE CLAIM MACROS
%
%  Almost nothing in this subject can be measured, which is exactly why the
%  kind of every claim is marked on the page. A reader can then see which
%  sentences are evidence and which are taste, and disagree with the second
%  kind without distrusting the first.
%
%  \reported{text}{sourcekey}
%      What the market does. TWO mandatory arguments, so a reported claim
%      with no source behind it is an "argument missing" error at build time
%      rather than something tools/check_sources.py has to catch afterwards.
%      Making the wrong form impossible beats making it detectable.
%
%  \reported[single]{text}{sourcekey}
%      The same, for a source graded single-source in sources.yaml. The
%      optional argument prints "one source" in the sentence, so the
%      qualification reaches the reader rather than living in a YAML file
%      they will never open. check_sources.py requires the optional argument
%      when the grade demands it AND refuses it when the grade does not, so
%      the marker cannot be sprinkled on for safety.
%
%  \judgement{text}
%      The author's opinion. Never set as fact, and never carrying a source,
%      because a citation on an opinion is a borrowed authority.
% ---------------------------------------------------------------------------
% A source mark prints a NUMBER, not the key.
%
% It printed the key at first, and keys are descriptive: several run to thirty
% characters, a superscript is one unbreakable run, and two of them in a row
% produced a 193pt overfull line against a 15pt budget. Measured, not guessed.
% Numbers come from tools/gen_sources.py, which also generates the register
% Appendix D prints -- so the list and the numbering cannot come apart.
%
% An unknown key prints loudly rather than silently. A citation that renders as
% nothing is a claim that has quietly lost its evidence.
\newcommand{\kitsourcenum}[2]{\expandafter\gdef\csname kitsrc@#1\endcsname{#2}}
\newcommand{\kitsourcemark}[1]{%
  \textsuperscript{\textcolor{kitsaid}{\scriptsize
    \ifcsname kitsrc@#1\endcsname
      \csname kitsrc@#1\endcsname
    \else
      \textcolor{kitwarn}{??#1}%
    \fi}}}

\newcommand{\source}[1]{\kitsourcemark{#1}}

\NewDocumentCommand{\reported}{o m m}{%
  \IfNoValueTF{#1}{%
    #2\kitsourcemark{#3}%
  }{%
    #2\ \textcolor{kitgrey}{\small(one source)}\kitsourcemark{#3}%
  }%
}

\newcommand{\judgement}[1]{%
  \textcolor{kitmine}{\textbf{\small$\diamond$}}\,\textit{#1}}

% ---------------------------------------------------------------------------
%  Drills
%
%  A drill is something the reader does, not something they read. It is a
%  distinct treatment because the whole argument of this book is that reading
%  a preparation guide is not preparation.
% ---------------------------------------------------------------------------
\newcounter{kitdrill}[chapter]
\renewcommand{\thekitdrill}{\thechapter.\arabic{kitdrill}}

\newtcolorbox[use counter=kitdrill]{drill}[1][]{%
  kit outlined=kitsaid,
  title={Drill \thekitdrill},
  #1}

% ---------------------------------------------------------------------------
%  A bank question
%
%  Six fixed parts, in this order, every time. The order is the content: it
%  forces the answer to separate what is asked from what is tested, and to
%  name the weaker answer rather than only the better one. A question missing
%  one of the six stops being comparable with the others, and
%  tools/check_structure.py refuses it.
% ---------------------------------------------------------------------------
\newcommand{\bankpart}[1]{%
  \par\needspace{2\baselineskip}%
  \vspace{4pt}\noindent{\bfseries\small\textcolor{kitink}{#1}}\par\nopagebreak\vspace{1pt}}

\newenvironment{bankq}[1]
  {\needspace{5\baselineskip}%
   \par\vspace{6pt}%
   {\color{kitink}\rule{\linewidth}{0.7pt}}\par\vspace{3pt}%
   \noindent{\large\bfseries #1}\par\vspace{2pt}}
  {\par\vspace{7pt}}

\newcommand{\qasked}{\bankpart{As asked}}
\newcommand{\qtested}{\bankpart{What is being tested}}
\newcommand{\qstaff}{\bankpart{The Staff answer}}
\newcommand{\qsenior}{\bankpart{The Senior answer that falls short}}
\newcommand{\qcontested}{\bankpart{Where this is contested}}
\newcommand{\qdrill}{\bankpart{Drill}}


% ---------------------------------------------------------------------------
%  Computed values
%
%  Appendix E prints ledgers about this repository -- how many chapters, how
%  many drills, how many sources at each grade. Those are facts about the tree
%  and they change with every pass, so they are NOT typed. tools/ledger.py
%  measures the tree and writes ledger.tex; the appendix reads it with \val{}.
%
%  The point is not convenience. A count typed into prose is a claim nobody
%  maintains, and this book spends a whole appendix refusing other people's
%  unmaintained counts. Printing its own from memory would be the same defect
%  with the author's name on it.
%
%  A missing key prints ?? loudly rather than silently, because a value that
%  quietly renders as nothing is a number that has stopped being true in the
%  one place nobody looks.
% ---------------------------------------------------------------------------
\newcommand{\kitledgerval}[2]{\expandafter\gdef\csname kitval@#1\endcsname{#2}}
\newcommand{\val}[1]{%
  \ifcsname kitval@#1\endcsname
    \csname kitval@#1\endcsname
  \else
    \textbf{\textcolor{kitwarn}{??#1??}}%
  \fi}

% ---------------------------------------------------------------------------
%  Small conveniences
% ---------------------------------------------------------------------------
% \code{} has to be able to break.
%
% This book names files constantly, and a path like templates/project-deep-dive.md
% is thirty characters of \texttt with no break opportunity: TeX either sets a
% 60pt overfull line or a very loose one. Measured before this fix: the longest
% such run produced a 66pt box.
%
% url's breaking machinery does the job -- it breaks after / . - and _ -- and
% it is configured here to prefer those points rather than to hyphenate, so a
% filename never acquires a hyphen that is not in it. Unlike \nolinkurl this
% takes an ordinary argument, so \code{\_} and \code{$\ldots$} still work.
\DeclareUrlCommand{\kiturl}{\urlstyle{tt}}
\newcommand{\code}[1]{\kiturl{#1}}
\newcommand{\term}[1]{\textit{#1}}
\newcommand{\dash}{\,---\,}

% A table that fills the measure, and that can break across a page.
%
% Two traps here, both measured rather than anticipated.
%
% \NewEnviron rather than \newenvironment: tabularx reads its body by scanning
% the input for a literal \end{tabularx}. Wrapped in an ordinary
% \newenvironment the source only contains \end{kittable}, so tabularx scans to
% the end of the file and the run dies with "File ended while scanning use of
% \TX@get@body" -- an error naming neither the table nor the chapter.
%
% xltabular rather than tabularx inside a center: a tabularx cannot break, so
% a comparison table taller than the room left produced an overfull vbox of
% 828pt -- a whole page of content with nowhere to go, reported by TeX as a
% page that came out too tall and naming no file. xltabular is a longtable, so
% it splits, and the class of defect disappears rather than being tuned.
\NewEnviron{kittable}[1]{%
  \par\vspace{4pt}%
  \begingroup\small
  \setlength{\LTpre}{0pt}\setlength{\LTpost}{0pt}%
  \begin{xltabular}{\linewidth}{#1}\toprule
  \BODY
  \bottomrule\end{xltabular}%
  \endgroup\vspace{2pt}}
